\section{Excluding calibrated aspect paths}\label{sec:aspect}

We complete the proof of Theorem~\ref{geo:main}. Suppose, to the contrary, that
\begin{equation}\label{asp:contradiction}
q_*<1+2s-10\kappa.
\end{equation}
The preceding construction supplies a calibrated entropy profile and finite
parameters
\[
a\ge0,\qquad b\le1,\qquad \lambda=a-b.
\]
We may take further simultaneous tangents at generic points of the calibrated
line. The support inequality, its equality on that line, the conditional
parameter profiles, and the predictability statement survive these operations
by Lemmas~\ref{ent:predictability} and~\ref{ent:tangent-realization}.
All actual entropy profiles below have been regularized as in
Lemma~\ref{ent:regularization}.
Thus a ``dense'' restriction means a restriction of subpower relative mass,
and support exponents used in a conditional block have the corresponding
typical-bin weight bounds. These conventions also apply to iterated tangents.

The cases below exhaust the finite parameters supplied by the
calibration.  The table records their roles; every comparison is
proved in the corresponding subsection.
\begin{center}
\small
\begin{tabular}{@{}p{.23\linewidth}p{.30\linewidth}p{.40\linewidth}@{}}
\toprule
Aspect & Further distinction & Comparison mechanism \\
\midrule
$\lambda=0$ & $a>0$ or $a=0$ & Isotropic profile, radial test,
and strip-axis reduction. \\
$\lambda>0$, $a>0$ & $\theta=0$; or $\theta>0$ with
$\lambda\ne1$ or $\lambda=1$ & Longitudinal comparisons,
pure-order costs, and the joint-coordinate tie. \\
$\lambda>0$, $a=0$ & Zero or positive full $C$-slope; for
$\theta>0$, causal interior or boundary & Static-coordinate separation and
a local constrained comparison. \\
$\lambda<0$ & $a>0$ or $a=0$ & Spatial axis prediction,
constrained pure orders, and longitudinal comparisons. \\
\bottomrule
\end{tabular}
\end{center}
Here $\theta$ is the additional conditional axis rate in
\eqref{eq:source-17}, not a spatial slope.  The backward gain
identity and reusable scalar comparisons are collected after the
stationary case.  All projection bounds below are applied to actual
realizing profiles; pure envelopes provide only the ordered costs
and overlap identities of Lemma~\ref{ent:signed-envelope}.

\subsection{Stationary aspect}

\begin{proposition}\label{asp:stationary}
There is no calibrated profile satisfying \eqref{asp:contradiction} with
\(\lambda=0\).
\end{proposition}

\begin{proof}
Here \(0\le a=b\le1\), and we use the isotropic profile
\(F(u,z,Y)\) of \eqref{eq:source-13}, on \(z\le u+Y\).
The reference-axis label is retained with the exact entropy
cancellation \eqref{ent:stationary-axis-cancellation}; the later
choice of a strip axis will also retain its earlier parent.
Suppose first that \(a>0\). Take a generic tangent on the calibrated line and
apply Lemma~\ref{ent:signed-envelope}, with \(u\) the only negative coordinate.
Its pure profile is \(Pu\), where \(P\) is constant and
\(\partial_uF\ge P\) throughout the domain. For \(T>0\), run backwards
from \((0,0,0)\) to \((T,0,-T)\), keeping \(z=0\), with the frozen
velocity cutoff in the upper envelope sufficiently fine for this
segment. The remaining time rate is at least \(s\), so the support
inequality and upper envelope give
\[
(2-q_*)T\le F(T,0,-T)\le(P-s)T.
\]
Consequently
\[
q_*\ge s+2-P,
\qquad\text{and hence}\qquad P>1-s.
\]

On the early side of the calibrated line, namely
\(u+aY<0\), \(z=u+Y\), the support inequality and equality on the line imply
an average bound \(\partial_uF+\partial_zF\le2\) when integrating toward
the line at fixed \(Y\). Since \(\partial_uF\ge P>1-s\), there is a generic
aligned point on this side with position rate
\(E:=\partial_zF<1+s\). Lemma~\ref{ent:peeling} gives an affine tangent in
the three isotropic coordinates there; denote its time rate by
\(\tau\ge s\). Simultaneously retain a tangent at the calibrated point with
the same velocity depth \(u\), whose time depth is \(-u/a\). This is later
than the early point and has finer position depth. The two sites have the
same physical velocity units, and the later profile still satisfies
\(\partial_uF\ge P\).

If \(a=0\), first peel off the constant coordinate on the line
\(u=0,z=Y\). In a tangent,
\[
F(u,z,Y)=J(u)+H_0(z,Y),
\]
where the second factor is defined on the full projected \((z,Y)\)-plane.
At a further generic point on the line, the two-coordinate order rule makes
\(H_0\) affine on each order half-plane. On \(z\ge Y\), write
\[
H_0(z,Y)=Ez+\tau Y.
\]
Then \(\tau\ge s\) and \(E+\tau=q_*\), so \(E<1+s\).
At \(Y=0,u=z>0\), the support inequality gives
\begin{equation}\label{asp:J-prefix}
J(u)\ge(2-E)u.
\end{equation}
The following projection block is based at the origin in this case.
Choose any fixed point on \(u=0,z=Y\) with \(Y>0\) as the later site;
the tests near it will use velocity depths near zero.

We give the common block argument. Write \(m\) for its base and let
\(b_0\in\mathbb R^2\) be normalized velocity. Let
\((d,c_0)\in\mathbb R\times\mathbb R^2\) be normalized time and position
in the fine-time frame, after subtracting known translations. Position in the
parent frame is therefore
\[
A_{b_0}(d,c_0):=c_0-d b_0
\]
to mesh accuracy. The profile and its regularization give the following
properties inside typical parents.
\begin{enumerate}[label=(\roman*)]
\item Velocity has a Frostman ball bound of some exponent strictly greater
than \(1-s\). For \(a>0\) this follows from \(\partial_uF\ge P\); for
\(a=0\) it follows from \eqref{asp:J-prefix} at every tested prefix depth.
\item Velocity cells and time-position cells have dense product incidence.
Indeed the refinements \((u,z,Y)=(1,0,0)\), \((0,1,1)\), and
\((1,1,1)\) have additive joint exponents. The time-position exponent is
\(E+\tau\).
\item Given a fine velocity bin \(b_0\), the image under \(A_{b_0}\) has
a Frostman \(E\) bound and at most \(m^{E+o(1)}\) bins. The fiber in one
image bin has at most \(m^{\tau+o(1)}\) time-position cells. These follow
by testing \(u=1\), \(0\le z\le1\), \(Y=0\), and the joint refinement,
rebasing with the known velocity.
\end{enumerate}

Suppose that every strip of power-small width has power-small velocity mass.
By Lemma~\ref{proj:radial}, applied with a slightly smaller ball exponent
at most one, directions \(b-b_0\) have an angular Frostman exponent
\(\sigma>1-s\) outside removable exceptional pairs. Product density and
Cauchy--Schwarz give a pin \(b_0\) and a dense set of such \(b\)'s that
share a dense set of time-position cells with \(b_0\). We may also arrange
\(|b-b_0|\ge m^{-o(1)}\).

Fix one of these overlaps. Keep its rich \(A_b\)-bins, each containing at
least \(m^{\tau-o(1)}\) cells. The support upper bound for \(A_b\) shows
that these retain full overlap exponent \(E+\tau\). Inside a fixed
\(A_b\)-bin the different cells give, up to bounded mesh multiplicity,
different \(d\)-bins: once \(d\) and \(A_b\) are given, \(c_0=A_b+db\)
is known to the tested accuracy. Moreover
\[
A_{b_0}=A_b+(b-b_0)d.
\]
Since \(|b-b_0|\ge m^{-o(1)}\), the rich bin therefore yields at least
\(m^{\tau-o(1)}\) distinct \(A_{b_0}\)-bins in one thick row parallel
to \(b-b_0\). Select perpendicular row labels separated by a sufficiently
large mesh constant. Their collections of point bins are disjoint up to
bounded multiplicity. There are at most \(m^{E+o(1)}\) point bins in all,
so the perpendicular projection of the retained \(A_{b_0}\)'s has at most
\(m^{E-\tau+o(1)}\) bins.

This restriction is dense for the Frostman image law, not merely large in
cardinality. Indeed full overlap exponent \(E+\tau\), together with the
fiber upper bound \(\tau\), gives \(E-o(1)\) image exponent. The typical
image bins have pre-trim weights at least \(m^{-E-o(1)}\), by
Lemma~\ref{ent:regularization}. Thus the image restriction has subpower
mass. We can consequently apply Lemma~\ref{proj:projection} to these
direction-dependent restrictions. It is contradicted by
\[
E-\tau<E,\qquad E-\tau<1,\qquad
E-\tau<(E+\sigma)/2.
\]
The first inequality uses \(\tau\ge s>0\), the second uses
\(E<1+s\), and the third follows from
\(E-2\tau<1-s<\sigma\). If \(E-\tau<0\), the support upper bound itself
is already impossible.

It follows that, in a dense family of parents, a dense amount of velocity
mass lies in strips of width \(m^{-\chi}\) for some fixed \(\chi>0\),
after subselection. Here the choice need not use one exceptionally rare
parent. To see this, average the maximal conditional strip mass over the
parents, using finite grids and bounded strip enlargements. If no fixed
width power had subpower average along a subsequence, each such average
would be power-small. Markov's inequality and diagonal selection over
countably many width powers would produce typical parents with all strips
power-light, contrary to the block argument. We retain the selected strip
mass across the parents, choosing its axis separately in each parent.

Choose a fixed neighborhood of the later site in which the time and
position cutoffs remain finer than those of the earlier parent. Use the
resulting parallel and normal coordinates there, and include the earlier
parent label. At later isotropic states
with positive velocity refinement, this label is recoverable up to bounded
ambiguity. Both time and position are finer than at the earlier parent;
rebasing position backwards multiplies velocity uncertainty by at most the
earlier time resolution. The label includes the earlier time bin, spatial
bins, and any fixed original slope label. Once it is known, the strip axis
is fixed, and switching isotropic grids costs only bounded multiplicity.
Thus the later isotropic profile and calibration are preserved under the
dense restriction.

The conditional time rate is also preserved. Given the base label, the
fixed original slope label, and a later time prefix, the earlier parent
time and its spatial bins, hence its selected axis, are already determined.
Conditioning on these extras therefore consumes no further time increment.
The linear frame change in velocity is used for position in the matched
units as well: position units are velocity units times the corresponding
time unit. In the two-site construction only the time and position units
change between the sites. This retains the scalar shear relation.
Lemma~\ref{ent:tangent-realization} allows these tests to be realized
simultaneously, with errors negligible in the final block units.

Now take mixed velocity depths \(l_1,l_2\in(0,\chi)\), with position
depths \(l_i+Y\). Given the parent, the normal velocity rate is zero
throughout this depth window. The parallel velocity rate is at least \(P\)
if \(a>0\), and at least \(J'(l_1)\) almost everywhere if \(a=0\).
For completeness, this lower bound is valid also at the aligned faces used
in projection tests. Compare a short parallel increment with the same
increment after making time and both position cutoffs finer, with the two
position cutoffs equal and with strict slack. At the finer state change
the normal velocity cutoff to the parallel cutoff at both endpoints; this
costs zero while both depths remain in \((0,\chi)\). The resulting
isotropic increment has the asserted lower rate. Dropping the extra
conditioning can only increase that increment by submodularity. Valid
interior comparisons pass to the aligned tangents by continuity.

Whenever this parallel rate exceeds \(1-s\), its aligned position rate
equals one by the scalar projection test \eqref{eq:source-16}, now in the
fixed axes. The distinct velocity, position, and time forms can be peeled
at generic mixed aligned states, so these are full partial rates there.

For \(a>0\), start on the later calibrated path with equal velocity
depths \(l_1=l_2=l\in(0,\chi)\). Choose the backward duration small
enough to stay in the chosen neighborhood and to keep both depths in
\((0,\chi)\). Hold the parallel depth \(l_1\) fixed and increase the
normal depth \(l_2\) at speed one. Then the normal position depth
\(l_2+Y\) is fixed, while the parallel position depth decreases at
speed one. The normal velocity has zero rate, so time contributes at
least \(s\) and parallel position contributes one. The advancing normal
depth is the larger width depth and therefore has coefficient
\(1-\kappa\) in the ellipse support term. Thus the profile minus
the ellipse width term decreases backwards at rate at least
\begin{equation}\label{asp:stationary-gain}
2+s-\kappa.
\end{equation}
To integrate generic derivatives, perturb the isotropic starting depth and
starting time slightly. The resulting segments sample all three mixed
alignment parameters. Integrate on generic segments and pass back by
Lipschitz continuity. Parent recoverability holds throughout because the
starting velocity depths remain inside the strip window. The initial
equality and support now contradict \eqref{asp:contradiction}.

For \(a=0\), \eqref{asp:J-prefix} supplies arbitrarily small positive
\(l\) with \(J'(l)>1-s\), on a set of positive measure. Fix a positive
comparison duration within the chosen later neighborhood and smaller
than \(\chi/2\). For each sufficiently small such generic \(l\), set
both initial depths to \(l\) and use generic later starting times.
The error of the initial isotropic state against the calibrated value
is \(O(l)\), and
including the parent adds no error in exponent. Sending \(l\) to zero
with the duration fixed again contradicts \eqref{asp:contradiction}, since
\(2+s-\kappa>1+2s-10\kappa\) for \(s\le1\).
\end{proof}

\subsection{Gain and scalar backward comparisons}

Henceforth \(\lambda\ne0\). Set
\[
u_*=1-s,\qquad h=(1+s)/2,\qquad
k_C=1+\kappa,\qquad k_B=1-\kappa.
\]
Testing states are fully aligned unless indicated otherwise:
\[
R=B-C,\qquad D=C+Y,\qquad A=B+Y.
\]
Put \(\mathcal H=F-k_CC-k_BB\), and write \(z(y)\) for a fully aligned
state with \(Y=y\). A backward segment of duration \(T\) starts at
\(z(0)\) and ends at \(z(-T)\); its \emph{gain} is
\(\mathcal H(z(0))-\mathcal H(z(-T))\).
Prescribe velocities by \(C'=-v_1\), \(B'=-v_2\), where primes denote
forward \(y\)-derivatives. Recall that
\[
(p,n_1,e,g,t,r)
=(\partial_CF,\partial_BF,\partial_DF,\partial_AF,
  \partial_YF,\partial_RF).
\]
The forward \(y\)-derivative of \(\mathcal H(z(y))\) is
\begin{equation}\label{eq:source-19}
t+(v_1-v_2)r-v_1p-v_2n_1
 +(1-v_1)e+(1-v_2)g+k_Cv_1+k_Bv_2.
\end{equation}
Indeed \(R'=v_1-v_2\), \(D'=1-v_1\), and \(A'=1-v_2\).
Integrating this derivative over \([-T,0]\) gives the gain. With the
increasing backward clock \(\tau=-y\), the potential in
Lemma~\ref{ent:control} is \(-\mathcal H(z(-\tau))\).
At a calibrated start, referenced to zero, \eqref{eq:source-11} gives
\(\mathcal H(z(-T))\ge-q_*T\), so the gain is at most \(q_*T\).
If an upper envelope \(U\ge F\) from \eqref{eq:source-14} agrees with
\(F\) at the start, its endpoint gain is a lower bound for the actual
gain. An initial error tending to zero is harmless.

We make precise how rates on generic sheets will be used. A singleton
rate obtained by peeling is a constant full partial in that tangent; a
group rate is the derivative along its specified diagonal. To integrate
an inequality valid on a sheet-open region, first integrate on almost all
slightly translated parallel segments and then pass to the chosen segment
by Lipschitz continuity. This also transfers an inequality to a boundary
approached from the indicated side. Subsequent off-sheet tests of a
singleton-affine tangent retain that singleton rate.

In particular, let \(\Gamma\) be the fully aligned plane
\(R=\lambda Y\) when \(\lambda>0\). For \(\lambda\ne1\), its spatial
partials are singleton constants in a generic first tangent, and the
parameter diagonal \((Y,R)\mapsto(Y+z,R+\lambda z)\) has rate at least
\(s\). The latter lower bound can be tested with spatial cutoffs fixed
in the interior of the separated tangent. Projection tests on its two
sides supply the subscripts \(s\) and \(\theta\) of
\eqref{eq:source-17} simultaneously. When \(\lambda=1\), the corresponding
grouped conclusion is \eqref{eq:source-18}. These facts hold in tangents
based anywhere on the tie, not only on the equality line.

\begin{lemma}[Scalar comparison]\label{asp:scalar}
Suppose a local restriction of an inherited profile has grouped coordinates
\(X,y,Z\), with domain \(Z\le X+y\), and no other local constraints.
The groups consist of disjoint nonnegative coordinate increments of the
underlying profile; additional fixed coordinates are allowed. At generic
aligned points write their rates as \(d_X,d_y,d_Z\). Suppose, also in the
tangent tests at the path points under consideration, that
\begin{equation}\label{eq:source-20}
d_y\ge t_0,
\qquad d_X>u_*/2\ \Longrightarrow\ d_Z\ge h.
\end{equation}
For \(k\in\{k_C,k_B\}\), the gain of this scalar profile minus \(kX\)
can reach rate \(t_0+h-\kappa\), with arbitrarily small further loss.
More precisely, along every path on \(Z=X+y\), parametrized by \(y\)
and with \(-X'(y)\in[\rho,1-\rho]\), almost every point admits arbitrarily short
straight backward comparisons with this rate, up to a loss tending to zero
with \(\rho\), and with velocity in the same interval.
\end{lemma}

\begin{proof}
Apply Lemma~\ref{ent:signed-envelope} separately along the path under
consideration. The negative coordinate is \(X\); \(y,Z\) strictly
increase. Slack is obtained by shifting \(y\) more than \(Z\).
At almost every path point the resulting tangent has pure slope \(p_0\),
and \(d_X\ge p_0\). If \(p_0>u_*/2\), velocity zero gives rate at least
\(t_0+h\) by \eqref{eq:source-20}.

If \(p_0\le u_*/2\), use velocity one. In centered offsets its ray is
\(X=-y,Z=0,y<0\). If \(m\in[\rho,1-\rho]\) is the original path speed,
choose \(\beta\in(y/m,y)\). Then
\[
-y>-m\beta,\qquad y>\beta,\qquad0>(1-m)\beta.
\]
These inequalities give the overlap condition of
Lemma~\ref{ent:signed-envelope}, locally up to alignment. Thus
\(d_X=p_0\) on the comparison ray, and integration gives rate at least
\[
t_0+k-p_0\ge t_0+1-\kappa-u_*/2=t_0+h-\kappa.
\]
Perturb velocities zero and one into the permitted interval. The loss is
arbitrarily small by the Lipschitz bounds. The strict comparison in the
tangent gives arbitrarily short comparisons in the original profile.
Any additional linear contribution to the gain is included in these same
endpoint calculations.
\end{proof}

\begin{lemma}[Longitudinal comparisons]\label{asp:longitudinal}
The following two comparisons are available locally, and in inherited
path tangents, under the stated derivative and conditional-time bounds.
\begin{enumerate}[label=(\roman*)]
\item \emph{Low longitudinal rate.} Fix \(v_1=1\), so \(D=d_0'\) is
constant. Suppose the backward cost of the varying \(C\)-term has been
bounded by \(P\le u_*/2\) and removed. This may be done in the pure upper
envelope when \(B\) is not negative, or where \(\partial_CF=P\) exactly.
If \(t\ge s\) in the local region, the total gain can reach
\(1+2s-\kappa\), with arbitrarily small further loss.
\item \emph{High longitudinal rate.} Fix \(v_1=0\), so \(C=c_0'\) is
constant. Suppose \(p\ge P>u_*/2\) in neighborhoods of the aligned states
being visited, up to the domain constraints, and \(t\ge s\). The same
gain \(1+2s-\kappa\) is available.
\end{enumerate}
\end{lemma}

\begin{proof}
For the low comparison use the scalar groups
\[
X=B,\qquad y=Y,\qquad Z=A=R+d_0'.
\]
For the upper envelope \eqref{eq:source-14}, the remaining function is
\[
F^+(Q)=F(n(-H_0),Q)-G(n(-H_0)).
\]
Its rates are those of the actual frozen-cutoff slice
\(Q\mapsto F(n(-H_0),Q)\), since the subtracted term is constant.
Choose \(H_0\) sufficiently large for all the following compact tests;
in particular \(C\) is fixed sufficiently fine in this slice.
Any additional negative \(D\) is held fixed in the pure part;
in the remaining part it is either \(d_0'\) or fixed finer, with
\(D<C+Y\). In the exact-slope version \(F-PC\) is locally independent
of \(C\). Freeze \(C\) at a slightly finer value to obtain strict slack
for \(B<C+R\) and \(D<C+Y\); on alignment this auxiliary slice agrees
with the remainder up to a constant.

To verify \eqref{eq:source-20}, retain \(R,A\) separately with
\(A\le D+R\), \(A\le B+Y\), while the other faces have slack. At a
generic scalar aligned point in any path tangent, peel off the singleton
rates for \(B,Y\); the \((R,A)\)-summand has affine diagonal rate.
Increase \(R\) separately to open \(A<D+R\), retaining \(A=B+Y\).
The time-projection block for \((B,Y,A)\) gives
\(g\ge f_s(n_1)\). Its fixed \(D\)-measurement rebases under time
refinement with bounded ambiguity because \(C\) is fine. The two marginal
laws and their dense product relation are the comparable-cutoff laws of
Lemma~\ref{ent:regularization}; thus this is the actual sheared block
test. If \(d_X=n_1>u_*/2\), it follows that \(g\ge h\).
The direction increasing \(B,R,A\) together preserves \(A=B+Y\)
and has rate \(n_1+r+g\ge n_1+h\). Integrate this inequality on
the displaced matching slices and pass back by continuity to the
\((R,A)\)-grouped slice. There the same direction has rate
\(n_1+d_Z\). The peeled singleton \(n_1\) is constant on both slices,
so subtraction gives \(d_Z\ge h\). The conditional parameter
increment with \(C,D\) fixed gives \(d_y\ge s\).

These constructions also apply after the signed-path tangent: take
simultaneous tangents of the indicated auxiliary slices, which have actual
realizing states. When freezing \(C\) is only local, cover the open path
portion by countably many neighborhoods with a single valid frozen value
on each. Lemma~\ref{asp:scalar} with \(t_0=s\) now gives total rate at least
\[
s+h-\kappa+k_C-P\ge1+2s-\kappa.
\]

For the high comparison the scalar groups are
\[
X=B=R+c_0',\qquad y=Y=D-c_0',\qquad Z=A.
\]
At full-aligned generic states, \(e\ge f_s(p)\ge h\). Hence the
\(y\)-group has rate at least \(s+h\). If \(g<h\), the projection
inequalities imply \(n_1\le u_*/2\). They also force \(r=0\): otherwise
\(g\ge f_r(e)\ge h\), a contradiction.

We justify transfer to the grouped slice, including its path tangents.
At a generic scalar aligned point, first peel off the fixed coordinate
\(C\) by a direction increasing every other coordinate; then separate
the \(X,y,Z\) groups. Their distinct forms are positive on a common
vector. The \(A\)-rate is therefore a constant singleton in this first
tangent. Test full-aligned generic points on nearby translated sheets,
allowing \(C\) to vary. The aligned \(y\)-direction also advances \(A\);
subtracting its constant rate transfers \(d_y=t+e\ge s+h\).
If this singleton rate is less than \(h\), the aligned \(X\)-direction,
with the same subtraction, transfers
\(d_X=n_1+r\le u_*/2\). Thus \eqref{eq:source-20} holds with
\(t_0=s+h\), and Lemma~\ref{asp:scalar} gives
\(s+2h-\kappa=1+2s-\kappa\).
\end{proof}

Whenever a prescribed-length comparison is needed in a free region, apply
Lemma~\ref{ent:control} to these pathwise short steps, taking \(\rho\)
small. If the construction occurs inside a tangent, its endpoint comparison
is enough: the effective straight velocity lies in the convex velocity box,
and the comparison passes back to arbitrarily short original steps. A
linear closed half-space constraint passes back as well. Pure backward
cost bounds may be integrated with the remaining gain, or subtracted at
the final endpoint from an upper envelope that is exact at the start.
Choose the additional endpoint, clipping, and control losses so that,
together with the loss already incurred in the selected comparison,
the total is less than \(10\kappa\).

\subsection{Increasing aspect with positive longitudinal speed}

\begin{proposition}\label{asp:increasing-positive}
There is no calibrated profile satisfying \eqref{asp:contradiction} with
\(\lambda>0\) and \(a>0\).
\end{proposition}

\begin{proof}
Take the conditional parameter tangent \eqref{eq:source-17} and a further
simultaneous signed-envelope tangent on the calibrated line. The negative
coordinates are \(C\), together with \(B\) if \(b>0\) and \(D\) if
\(a>1\). Their pure domain is full: the parameters \(R,Y\) in mixed
constraints increase, and negative cutoffs at an earlier path time can be
paired with the other cutoffs at a later time. Large positive shifts in
\(R,Y\) open the required slack. On compact aligned tests all negatives can
be replaced by \(n(-H_0)\) in \eqref{eq:source-14} for sufficiently large
\(H_0\). When relevant, the differences \(C-B\) and \(C-D\) of the
shifted negatives increase strictly, since \(a>b\) and \(a>a-1\).

Let \(P\in[0,1]\) be the pure marginal rate of \(C\). If \(P=0\), test
the upper envelope backwards for unit time with \(v_1=M_0,v_2=0\),
where \(M_0\) is large and fixed. The \(C\)-cost is zero, \(B\) is
constant, and \(Y,A,R\) decrease backwards. The only possible positive
backward cost is that of \(D\), at most \(M_0-1\), whether it belongs
to the pure or remaining profile. The gain is at least
\[
k_CM_0-(M_0-1)=1+\kappa M_0,
\]
contradicting support for large \(M_0\). Comparisons along binding faces
follow by translating the segments into the relative interior first.
Henceforth \(P>0\).

We divide the remaining proof into three cases. When \(\theta=0\),
we use longitudinal comparisons, upper-envelope tests, and full-alignment
projection bounds. When \(\theta>0\) and \(\lambda\ne1\), we use the
projection inequalities on the tied plane and then separate the possible
velocities. When \(\theta>0\) and \(\lambda=1\), the coincident spatial
forms require the grouped projection bounds.

\subsubsection*{No additional conditional axis rate}

First suppose \(\theta=0\), so \(t\ge s\) everywhere. If \(b\le0\)
and \(P\le u_*/2\), apply the low comparison of
Lemma~\ref{asp:longitudinal} to the upper envelope starting at zero.
Its pure cost along \(v_1=1\) is at most \(P\), and any additional
negative \(D\) is constant. If \(b\le0\) and \(P>u_*/2\), use the
high comparison in the actual profile. On its backward paths from
\(C=0\), the pure \(C\)-rate is \(P\): if \(D\) is negative, its
value \(D=y<0\) is strictly coarser than \(C=0\) in normalized negative
units. Thus \(p\ge P\) locally at the visited states. Both comparisons
contradict support.

Suppose \(b>0\), and let \(Q_1\) be the pure marginal of \(B\).
If \(P+Q_1\le u_*\), the upper-envelope test \(v_1=v_2=1\) keeps
\(D,A,R\) fixed and costs at most \(P+Q_1\) in the pure part. The
remaining time rate is at least \(s\), so its gain is at least
\(2+s-(P+Q_1)\ge1+2s\). If \(Q_1=0\) and \(P>u_*\), the pure
profile is independent of \(B\), and the high comparison has \(p\ge P\)
locally at its visited states, with any negative \(D\) strictly coarser.

The remaining possibility is \(P,Q_1>0\) and \(P+Q_1>u_*\).
Consider the aligned sheet
\begin{equation}\label{eq:source-21}
\begin{gathered}
C=-a\alpha,\qquad B=-b\alpha,\qquad R=B-C,\\
Y=y,\qquad D=C+y,\qquad A=B+y .
\end{gathered}
\end{equation}
Initially restrict to \(y<\alpha\). The free \(y\)-direction increases
precisely \(D,A,Y\), so at a generic point peeling splits off
\((C,B,R)\). The remaining three forms are distinct, since
\(a,b>0\) and \(a\ne b\). Therefore \(e,g,t\) are singleton constants
in a first tangent. Keeping \(Y\) fixed, use the two perturbations
\begin{equation}\label{asp:aligned-perturbations}
\begin{aligned}
(C,B,D,A,Y,R)&\longmapsto(C+\delta,B,D+\delta,A,Y,R-\delta),\\
(C,B,D,A,Y,R)&\longmapsto(C,B+\delta,D,A+\delta,Y,R+\delta),
\end{aligned}
\qquad \delta>0.
\end{equation}
They preserve full alignment and make \(C\), respectively \(B\),
strictly finest among these two normalized negative depths. Any negative
\(D\) is strictly coarser because \(y<\alpha\), and remains so for
sufficiently small \(\delta\). The pure lower bounds give
\(p\ge P\) on the first order side and \(n_1\ge Q_1\) on the
second. The full-aligned projection tests therefore give the following
bounds for the same singleton constants \(e,g,t\):
\[
e\ge f_s(P),\qquad g\ge f_s(Q_1),\qquad t\ge s.
\]
The two clipped position bounds sum to at least \(1+s\): if neither
clips their sum is \(2s+P+Q_1>1+s\); if one clips the other is
at least \(s\). At fixed \(\alpha\), the gain is at
least \(1+2s\). Integrate along generic sheet segments and let
\(\alpha\to0\). This gives the contradiction on the backward ray
\(v_1=v_2=0\), \(y<0\).

\subsubsection*{The tied plane for \(\theta>0\), \(\lambda\ne1\)}

Now assume \(\theta>0\), \(\lambda\ne1\). At generic points of
\(\Gamma\), the parameter diagonal rate is at least \(s\), and
\begin{equation}\label{eq:source-22}
e\ge f_s(p),\qquad g\ge f_s(n_1),\qquad
n_1\ge f_\theta(p),\qquad g\ge f_\theta(e).
\end{equation}
Set
\[
K_L=p+e,\qquad K_N=n_1+g,\qquad K_\Sigma=K_L+K_N.
\]
In particular \(K_N\ge K_L\) and \(g\ge e\). Parameterize \(\Gamma\)
by \(L_0=C+ay\) and \(y\), and average over thin strips on each side
of \(L_0=0\) in a fixed time interval. Integrating the support bound
from the equality line gives average \(K_\Sigma\le2\) on the minus
side and average \(K_\Sigma\ge2\) on the plus side. For the gain
on the tie with \(v_1=m,v_2=m-\lambda\), formula
\eqref{eq:source-19} becomes
\begin{equation}\label{eq:source-23}
\mathcal D_\Gamma(m)
=\text{parameter diagonal rate}+e+g+m(2-K_\Sigma)
 +\lambda(K_N-1+\kappa).
\end{equation}
The strip average of \(\mathcal D_\Gamma(a)\) tends to \(q_*\) on
either side as the strip shrinks, by calibration and Lipschitz continuity.
The plus-side average of \(K_N\) is at least one.

Write \(P,N_-,E_-\) for the minus-side pure rates of the negative
coordinates \(C,B,D\), using zero for an absent negative coordinate.
Here the \(C\)-rate is its marginal \(P\), by the ordering of
\(L_0\) divided by the negative speeds. These are lower bounds for the
actual partials. On the plus side let the pure rates be \(p_+,N_+,E_+\).
They are exact actual negative partials there. Indeed all negative
cutoffs are strictly finer than the path at time \(y\), while the
other cutoffs are finer or equal; replacing the latter by their path
values is valid. Increase all negatives strictly, with the \(C\)
increment at least as large as those of any negative \(B,D\).
This preserves the constraints at both tuples. Overlap gives equality
of the increment with its pure value. Each partial already has its
pure affine order slope as a lower bound; in the peeled tangent,
equality of this positive combination forces individual equality.
Interior derivative comparisons pass to these tangents also at alignment.

The weighted pure rate does not depend on the order:
\begin{equation}\label{eq:source-24}
J=aP+b_+N_-+(a-1)_+E_-
 =a p_++b_+N_++(a-1)_+E_+.
\end{equation}
The negative speeds in this use are distinct because \(\lambda\ne1\).
Marginal maximality gives \(p_+\le P\). The minus-side averages imply
\begin{equation}\label{eq:source-25}
P+\max(P+\theta,N_-)\le u_*,
\qquad P+E_-\le1.
\end{equation}
Here is the clipping argument. Since \(K_N\ge K_L\), the average of
\(K_L\) on the minus side is at most one. If \(P+s\ge1\), then
\(K_L\ge P+1>1\), a contradiction. Put
\(l=\max(f_\theta(P),N_-)>0\). Then
\[
K_\Sigma\ge P+(P+s)+l+f_s(l).
\]
If \(l+s\ge1\), this is at least \(2+2P>2\). Thus \(l+s<1\),
which also gives \(P+\theta<1\) and
\(l=\max(P+\theta,N_-)\). The average bound now gives
\(P+l\le1-s\). Finally \(K_L\ge P+E_-\) proves the second
inequality. We will use
\[
2P\le u_*,\qquad P+N_-\le u_*,
\qquad\theta\le u_*-2P.
\]
When \(s=1\), these already contradict \(P>0\). Otherwise the
following cases cover the present parameters.

\paragraph{Case \(a>1\), \(b\le0\).}
On the plus side, direct expansion gives
\[
\mathcal D_\Gamma(a)\ge s+a-J+g+b(1-K_N)+\kappa\lambda.
\]
The term with \(b\) is nonnegative after averaging, because the
plus-side average of \(K_N\) is at least one. Also
\[
J=aP+(a-1)E_-\le a-1+P.
\]
If \(E_+\ge s+P\), then \(g\ge e=E_+\), giving average gain at
least \(1+2s+\kappa\lambda\). If \(E_+<s+P\), then
\(p_++E_+\le2P+s\le1\), and therefore
\[
J=(a-1)(p_++E_+)+p_+\le a-1+p_+.
\]
When \(p_+>0\), \(g\ge e\ge f_s(p_+)=s+p_+\), again giving the
required gain. When \(p_+=0\), choose fixed \(M_0>a\) large and
use \(v_1=M_0,v_2=0\) in the upper envelope. The negative tuple
\((C,D)=(M_0,M_0-1)\) has the plus order, since
\[
\frac{M_0-1}{a-1}>\frac{M_0}{a}.
\]
Its pure cost is at most \(M_0-1\), and the varying remaining
coordinates decrease backwards. The gain is at least
\(1+\kappa M_0\), contradicting support for large \(M_0\).

\paragraph{Case \(a\ge1\), \(b>0\).}
Here \(0<b\le1\). The exact plus-side rates give
\[
\begin{split}
\mathcal D_\Gamma(a)
&\ge s+a+b-J+(1-b)g+\kappa\lambda\\
&\ge s+1-P+b(1-N_-)+(1-b)g+\kappa\lambda,
\end{split}
\]
using \(J\le a-1+P+bN_-\). We have \(1-N_-\ge s+P\).
If \(N_+\ge P\), then \(g\ge f_s(N_+)\ge s+P\) pointwise.
If \(N_+<P\), use \(p_+\le P\), \(e\le g\), and exactness to
get \(K_\Sigma\le2P+2g\). Since its plus-side average is at least
two, the average of \(g\) is at least \(1-P\ge s+P\).
In both alternatives the average gain is at least
\(1+2s+\kappa\lambda\). At \(a=1\), \(D\) is absent from the
negative group and its coefficient is zero, so the same displayed
calculation applies.

\paragraph{Case \(0<a\le1\), \(b\le0\), \(\lambda>1\).}
Only \(C\) is negative, and \(p=P\) exactly on the plus side.
Formula~\eqref{eq:source-23} gives
\[
\mathcal D_\Gamma(a)-\mathcal D_\Gamma(1)
=(a-1)(2-K_\Sigma),
\]
whose plus-side average is nonnegative. Moreover,
\[
\mathcal D_\Gamma(1)
\ge s+1-P+g+(\lambda-1)(K_N-1)+\kappa\lambda.
\]
Here \(g\ge e\ge s+P\), and the last \(K_N-1\) term is
nonnegative in average. This again gives average gain at least
\(1+2s+\kappa\lambda\).

\paragraph{Case \(0<b<a<1\).}
We have \(0<\lambda<1\), and the negatives are \(C,B\).
The plus-side rate \(N_+=Q_1\) is the \(B\)-marginal.
It is positive: if \(Q_1=0\), then \(N_-=0\), and
\eqref{eq:source-24} gives \(p_+=P>0\), contradicting
\(n_1=0\) and \(n_1\ge f_\theta(p)\) on the plus side.

If \(P+Q_1\ge u_*\), use \eqref{eq:source-21} with
\(\alpha<y\). This is strictly causal, since
\(R=\lambda\alpha<\lambda Y\). At a generic point its first
tangent has singleton constants \(e,g,t\). Apply the two perturbations
\eqref{asp:aligned-perturbations}, taking \(\delta<\lambda(y-\alpha)\)
to preserve the strict causal inequality. The two negative orders give
\[
e\ge s+P,\qquad g\ge f_s(Q_1),\qquad t\ge s;
\]
the first bound is unclipped because \(2P\le u_*\).
For increasing \(\alpha,y\) together, the gain rate is at least
\begin{equation}\label{eq:source-26}
a+b+\kappa\lambda+s-J
 +(1-a)(s+P)+(1-b)f_s(Q_1).
\end{equation}
To justify the cost \(J\) at the negative-order tie, perturb within
the tangent to \(B/b>C/a\), equivalently \(C+aR/\lambda>0\)
on alignment. At \(\beta=R/\lambda\) both negatives are
strictly finer than the calibrated path, and all positives are finer
or equal; in particular \(Y>\beta\), with the corresponding
position margins. Replacing positives by their path values is valid,
still with \(B=C+R\). Ordered increases of \(C,B\), with the
\(C\)-increment at least as large, preserve validity. Overlap gives
exact negative rates on this open order side, with weighted sum
\(J\); these comparisons pass to the tangent. The remaining axis
term \(\lambda r\) is nonnegative. Integrate the derivative bound
on perturbed parallel segments that retain strict order, then approach
the tie. This proves \eqref{eq:source-26}. Parallel sheets approaching
\(\alpha=y\) transfer it to the calibrated rate.

If \(Q_1\ge u_*\), subtract \(1+2s\) from
\eqref{eq:source-26}, using \(f_s(Q_1)=1\), to obtain
\[
\begin{split}
&\kappa\lambda+a(u_*-P)-bN_-+(1-a)P\\
&\qquad=\kappa\lambda+a(u_*-P-N_-)+(a-b)N_-+(1-a)P\ge0.
\end{split}
\]
If \(0<Q_1<u_*\), the difference is
\[
\begin{split}
&\kappa\lambda+P+Q_1-u_*+a(u_*-2P)+b(u_*-N_--Q_1)\\
&\quad=\kappa\lambda+(1-b)(P+Q_1-u_*)
 +(a-b)(u_*-2P)+b(u_*-P-N_-)\ge0.
\end{split}
\]
Each term has the stated sign by \eqref{eq:source-25}.

If instead \(P+Q_1\le u_*\), put
\(A_0=u_*-(p_++Q_1)\ge0\). The upper-envelope test
\(v_1=v_2=1\) has pure cost \(p_++Q_1\), because its
negative tuple \(C=B=1\) has the plus order. The remaining
time rate is at least \(s-\lambda\theta\). Its gain minus
\(1+2s\) is therefore at least
\[
A_0-\lambda\theta.
\]
For a second test follow \(\mathcal D_\Gamma(1)\) backwards
from zero. Its ray has \(L_0=(a-1)y>0\) for \(y<0\),
so \(p=p_+\), \(n_1=Q_1\) exactly, and
\(g\ge s+Q_1\), since \(0<Q_1<u_*\).
Formula~\eqref{eq:source-23} gives gain minus \(1+2s\)
at least
\[
A_0-\lambda(u_*-2Q_1)+\kappa\lambda.
\]
One test suffices. If \(Q_1\ge u_*/2\), the second expression
is nonnegative. Otherwise set \(d=u_*-2\max(P,Q_1)\ge0\).
Then
\[
A_0\ge u_*-P-Q_1\ge d,\qquad
0\le\min(\theta,u_*-2Q_1)\le d.
\]
Since \(\lambda<1\), the larger of the two lower bounds is
nonnegative even after discarding \(\kappa\lambda\).

\paragraph{Case \(0<a<1\), \(b\le0\), \(0<\lambda<1\).}
Only \(C\) is negative, with \(0<P\le u_*/2\).
Start at \(C=B=L>0,Y=0\) on alignment, where \(L\) is
arbitrarily small and generic, and run backwards for unit time with
\begin{equation}\label{eq:source-27}
C=L-y,\qquad B=L+x,\qquad
v_2=-x'\in[\rho,1-\rho],\qquad R=x+y\le\lambda y.
\end{equation}
Choose \(\rho\) small also relative to \(1-\lambda\).
These paths satisfy \(0\le x\le-y\). At \(y<0\), the
actual \(C\)-partial is \(P\) locally by overlap.
Indeed \(R\ge y\) and \(\lambda\ge a\) give
\(R/\lambda\ge y/a\). Choose
\(\beta\in((y-L)/a,y/a)\). Then \(C>-a\beta\), and
each of
\[
B=L+x,\quad D=L,\quad A=L+x+y,\quad Y=y,\quad R=x+y
\]
is strictly finer than its path value at \(\beta\).
For \(B\) this uses \(b\le0\); for \(D,A,Y\) it uses
\(a<1,b\le0,\beta<y/a\); for \(R\) it uses the
preceding inequality. The comparison is valid in a neighborhood
up to alignment.

In the strict causal interior \(t\ge s\), so the exact-slope
low comparison gives the pathwise short steps. On the boundary
\(R=\lambda y\), almost every \(y\) is a peeling point of
\(\Gamma\) for generic \(L\): the map
\((L,y)\mapsto(C,Y)=(L-y,y)\) is invertible, so Fubini
applies. This exceptional set in \(y\) is independent of the
path, since its boundary state is determined by \(L,y\).
In the first tangent there, \(p=P\) and \(g\ge s+P\).
Use \(v_1=1,v_2=0\) to enter the strict causal side
backwards. There \(t\ge s,r\ge0\), and the spatial rates
are constant throughout the tangent. The gain reaches
\[
s+k_C-P+g\ge1+2s+\kappa.
\]
Perturb \(v_2\) to \(\rho<1-\lambda\) to retain the
inward direction and the strict gain. Lemma~\ref{ent:control}
gives the endpoint contradiction, since the starting error tends
to zero with \(L\).

\subsubsection*{The coincident tie \(\lambda=1\)}

Finally suppose \(\lambda=1,\theta>0\), with \(P>0\).
Use paths on \(\Gamma\) through zero with
\(v_1=m\in[\rho,1-\rho]\), \(v_2=m-1\).
At \(y<0\) they have \(C>0\), \(D=B<0\).
Thus \(p\ge P\) locally: \(C\) is strictly finest if
\(D,B\) are also negative. Use the grouped domain
\[
B=D=Z_1\le C+Y,\qquad A=Z_2\le Z_1+Y,\qquad R=Y.
\]
Its parameter diagonal rate is at least \(s\).
The \(Z_1\)-diagonal rate \(H\) and the \(Z_2\)-rate
\(g\) satisfy \eqref{eq:source-18} at generic aligned
states and in the relevant inherited tangents. The slice retains
the binding alignment required by the tied projection argument.

Along each path, \(C\) alone decreases, while \(Y,Z_1,Z_2\)
strictly increase. Take its signed tangent, with positive shifts
of these latter coordinates opening both constraints, and let
\(p_0\) be the pure \(C\)-slope. Exactness in overlap
somewhere, together with the inherited \(p\ge P\), gives
\(p_0\ge P>0\). If \(p_0\le u_*\), test
\(v_1=1,v_2=0\) on the tie. Its ray is
\(C=-y,Z_1=0,Z_2=y\); a path time strictly between
\(y/m\) and \(y\) gives exact \(p=p_0\), since
the other grouped cutoffs are strictly finer. Its gain is at least
\[
s+k_C-p_0+g\ge s+k_C-p_0+(s+p_0)=1+2s+\kappa.
\]
If \(p_0>u_*\), use \(v_1=0,v_2=-1\).
The pure lower bound and \eqref{eq:source-18} give
\(H,g\ge1\), so the gain is at least
\[
s+H+2g-k_B\ge2+s+\kappa\ge1+2s+\kappa.
\]
Integrate along nearby generic tie segments and perturb the endpoint
values \(m=0,1\) into \([\rho,1-\rho]\), retaining
\(v_2=m-1\). The tie and velocity box are preserved.
These are the required pathwise control steps, giving the final
contradiction in this case.
\end{proof}

\begin{lemma}[Local constrained comparison]\label{asp:constrained}
Let an inherited full profile be based on the parameter tie, with
\(0<\lambda<1\), \(\theta>0\), the conditional profile
\eqref{eq:source-17}, and the inherited tie projection inequalities.
Suppose its full \(C\)-partial is a constant
\(P\in(0,u_*/2]\). Without requiring calibration of this profile,
there is a backward endpoint comparison from its origin, respecting
\(R\le\lambda Y\), whose gain per unit time is at least
\(1+2s-4\kappa\). Its velocity is in a box with \(v_1=1\)
and \(v_2\in[\rho,1-\rho]\), for sufficiently small \(\rho>0\).
\end{lemma}

\begin{proof}
Use the paths \eqref{eq:source-27} from a small generic
\(C=B=L,Y=0\). In the strict causal interior the constant
full \(C\)-slope and \(t\ge s\) supply the low longitudinal
comparison. At generic boundary points, the tie inequalities give
\(g\ge s+P\), and the spatial singleton rates in the first
tangent are constant. The backward direction \(v_1=1,v_2=0\)
enters the strict causal side and has gain at least
\(s+k_C-P+g\ge1+2s+\kappa\). Perturbing to
\(v_2=\rho<1-\lambda\) retains this strict gain.
The generic-\(L\) argument in Proposition~\ref{asp:increasing-positive}
gives these steps almost everywhere along every admissible path,
including its boundary epochs. Lemma~\ref{ent:control} therefore
gives a fixed-duration endpoint comparison, with arbitrarily small
additional loss.

Translate both endpoints by the same aligned depth shift to move
the start from \(L\) to zero. This changes the gain by \(O(L)\);
it does not alter \(R,Y\), so it preserves the half-space.
Choose \(L\) and the further clipping and control losses so that,
together with the \(\kappa\)-loss in the low longitudinal comparison,
the total loss per unit time is less than \(4\kappa\). The straight endpoint
displacement lies in the convex velocity box and respects the
half-space. It consequently certifies a short step upon passing
back from a tangent.
\end{proof}

\subsection{Increasing aspect with zero longitudinal speed}

\begin{proposition}\label{asp:increasing-zero}
There is no calibrated profile satisfying \eqref{asp:contradiction}
with \(\lambda>0\) and \(a=0\).
\end{proposition}

\begin{proof}
Here \(b=-\lambda\). Take \eqref{eq:source-17} and peel off
the constant coordinate \(C\) in a generic calibrated tangent;
all other coordinates strictly increase. Its additive contribution
is a function of \(C\) alone. Use aligned backward paths through
zero with
\[
\rho\le v_1\le M_0,\qquad -(\lambda+1)\le v_2\le1,
\]
where \(M_0\) is fixed sufficiently large relative to
\(\lambda\) and \(1/\kappa\). If \(\theta>0\), impose
\(R\le\lambda Y\). Along every such path, the \(C\)-summand
is differentiable almost everywhere because \(C\) moves
monotonically with speed at least \(\rho\). In a tangent
there its full slope \(p\) is constant. We verify the control
property in this tangent.

If \(p=0\), take \(v_1=M_0,v_2=0\).
After removing \(C\), the only backward cost is at most
\(M_0-1\), from \(D\), and the gain is at least
\(1+\kappa M_0\). Taking \(M_0\ge\lambda\) also
preserves the causal half-space at its boundary.
Suppose \(p>0\). If \(\theta=0\), or at a strictly
causal state, \(t\ge s\) throughout a neighborhood and its
tangents. Use the low longitudinal comparison for \(p\le u_*/2\)
and the high comparison for \(p>u_*/2\), perturbing
\(v_1=0\) to \(\rho\) in the latter.

It remains to consider boundary states with \(\theta>0\). The path
point need not be generic on \(\Gamma\). Its tangent nevertheless has
constant full \(C\)-slope \(p\). Each of the first three cases below
chooses a fixed tie direction using only \(p,\lambda\). Its gain bound
holds at generic points of the tangent's tie; integrate on nearby
generic parallel tie segments and pass by continuity to the segment
from the tangent origin. This gives the endpoint comparison required
by Lemma~\ref{ent:control}. The fourth case supplies such a comparison
directly.
\begin{enumerate}[label=(\roman*)]
\item For \(\lambda=1\), use \eqref{eq:source-18}.
If \(0<p\le u_*\), the tie test \(v_1=1,v_2=0\)
has gain at least \(s+k_C-p+(s+p)=1+2s+\kappa\).
If \(p>u_*\), the test \(v_1=0,v_2=-1\) has
gain at least \(s+1+2-k_B\). Perturb along the tie;
in the second test use \(v_1=\rho,v_2=\rho-1\).
\item For \(\lambda\ne1\), \(p>u_*/2\), the spatial
inequalities \eqref{eq:source-22} give
\(e,g\ge h\) and \(K_N\ge K_L\ge1\).
Formula~\eqref{eq:source-23} with \(m=0\) gives
at least \(s+2h=1+2s\). A small tie perturbation makes
\(v_1\ge\rho\).
\item For \(\lambda>1\), \(0<p\le u_*/2\), take
\(m=(\lambda+1)/2\) on the tie. Formula~\eqref{eq:source-23}
gives
\[
\mathcal D_\Gamma(m)\ge
s+1-p+g+\frac{\lambda-1}{2}(K_N-K_L)+\kappa\lambda.
\]
Here \(g\ge p+s\), since \(p+s\le h\le1\), and
the difference term is nonnegative. The rate is at least
\(1+2s\), and the velocities fit the chosen box.
\item For \(0<\lambda<1\), \(0<p\le u_*/2\),
use Lemma~\ref{asp:constrained}, with its constant full
\(C\)-slope \(P=p\). This enters or follows the causal
region and gives a strict endpoint comparison in the tangent.
\end{enumerate}
All these comparisons, including any secondary control construction
inside a tangent, retain a fixed positive margin over \(q_*\)
after the chosen losses. Their endpoint displacements lie in the
convex velocity box and respect the constant half-space constraint.
Lemma~\ref{ent:control} therefore contradicts
\eqref{eq:source-11} for the original paths.
\end{proof}

\subsection{Decreasing aspect}

\begin{proposition}\label{asp:decreasing}
There is no calibrated profile satisfying \eqref{asp:contradiction}
with \(\lambda<0\).
\end{proposition}

\begin{proof}
Write \(c=b-a=-\lambda>0\). Then \(0\le a<b\le1\).
We have \(t\ge s\) everywhere. The negative coordinates on the
calibrated line are \(B,R\), and \(C\) when \(a>0\).
Take the signed-envelope tangent with the slack shifts described in
Lemma~\ref{ent:signed-envelope}. In this case it is essential to
retain the additional prediction \eqref{eq:source-12}.

First suppose \(a>0\). The pure domain is \(B\le C+R\).
Let \(P\) be the constant pure \(C\)-partial on
\(C/a>\max\{B/b,R/c\}\), and let \(S\) be the constant pure
\(R\)-partial on \(R/c>\max\{C/a,B/b\}\).
Lemma~\ref{ent:signed-envelope} supplies these constants and makes each
a global upper bound for its corresponding pure partial.
On the aligned sector in which \(C\) is ahead,
\(C/a>R/c\) and \(B=C+R\), the pure profile has the form
\begin{equation}\label{eq:source-28}
G(C,B,R)=PC+Q_2B
\end{equation}
after referencing at zero, with constant \(Q_2\ge0\).
To prove this identity, note first that \(C/a>B/b\) as
well, since \(B/b=(a/b)(C/a)+(c/b)(R/c)\).
The pure \(C\)-rate is consequently \(P\).
The pure \(R\)-rate is zero there. Indeed in the external slices
used to construct the pure profile, \eqref{eq:source-12}
gives flatness in \(R\) when
\(C/a,B/b>R/c\); the other coordinates are fixed strictly
finer. This flatness holds on an open order region of the slice
tangents and passes to \(G\).

Thus varying \(C\) from the order tie at fixed \(B\), with
\(R=B-C\), has slope \(P\). On the tie,
\(C=(a/b)B\), the pure total is affine. Integrating from it
proves \eqref{eq:source-28}. Nonnegativity of the \(B\)-increment
gives \(Q_2\ge0\). The same \(B\)-rate holds locally in the
interior: the \(C,R\) partials there are constant, and one can
approach the aligned boundary keeping \(B\) fixed.
Thus the partials of \(G\) near this aligned sector are
\((P,Q_2,0)\) in the \((C,B,R)\) coordinates. The signed-envelope
lower comparison gives \(p\ge P\) and \(n_1\ge Q_2\) for the actual
profile.

The shift to \(n(-H_0)\) in \eqref{eq:source-14} is valid
on compact aligned tests because \(a,b,c>0\). If \(P=0\),
use \(v_1=M_0,v_2=0\) backwards from zero. The target pure
tuple is \(C=M_0,B=0,R=-M_0\), in the \(C\)-ahead sector,
so \eqref{eq:source-28} gives zero pure cost. The remaining
cost is at most \(M_0-1\), from \(D\); the other varying
remaining coordinates decrease backwards. Thus the gain is at
least \(1+\kappa M_0\), contradicting support for large fixed
\(M_0\).

If \(P+Q_2\le u_*\), use \(v_1=v_2=1\). Its pure tuple
\(C=B=1,R=0\) is again in the \(C\)-ahead sector and
costs \(P+Q_2\). The coordinates \(D,A\) are fixed,
while the remaining time rate is at least \(s\). The gain
is therefore at least \(2+s-(P+Q_2)\ge1+2s\).

If \(P,Q_2>0\) and \(P+Q_2>u_*\), use the sheet
\eqref{eq:source-21} with \(y<\alpha\), and the comparison
\(v_1=v_2=0\). At generic sheet points the same peeling
separates \(e,g,t\): the free \(y\)-direction first splits
off \(C,B,R\), after which the three remaining forms are
distinct. The first perturbation in \eqref{asp:aligned-perturbations}
enters the \(C\)-ahead sector while preserving full alignment.
Formula~\eqref{eq:source-28}, the pure lower bounds,
and the full-aligned projections give simultaneously
\[
e\ge f_s(P),\qquad g\ge f_s(Q_2),\qquad t\ge s.
\]
As above, the two positive input rates with sum greater than
\(1-s\) force \(e+g\ge1+s\). Integrate on generic
parallel segments and pass to \(\alpha=0\), obtaining gain
at least \(1+2s\).

The only remaining possibility with \(a>0\) is
\(P>u_*\), \(Q_2=0\). On the same sheet the
\(C\)-ahead perturbation gives \(e\ge1\).
If \(S>0\), the second perturbation in
\eqref{asp:aligned-perturbations} enters the \(R\)-ahead sector
\(R/c>C/a\), still with \(B=C+R\). There the actual
tilt rate is at least \(S\). Since \(e\) is a constant
singleton in the first tangent, the tilt projection inequality
\(g\ge f_r(e)\) gives \(g\ge1\). Transfer this to the
sheet and use \(v_1=v_2=0\), whose gain is at least
\(s+2\ge1+2s\).

If \(S=0\), the global upper bound for the pure \(R\)-partial makes \(G\)
independent of \(R\). It therefore projects to an ordinary
submodular profile on full \((C,B)\)-space: any \((C,B)\)
is feasible after taking \(R\) sufficiently fine, and the
value does not depend on that choice. If necessary take a further
simultaneous generic tangent on the calibrated line so that this
two-variable pure profile has ordinary order slopes, by the
full-domain order rule of Lemma~\ref{ent:signed-envelope}.
Its lower comparisons with the actual profile and
\eqref{eq:source-28} are preserved. If the resulting \(B\)
marginal is positive, the second perturbation in
\eqref{asp:aligned-perturbations} enters its \(B\)-ahead
order, which is the \(R\)-ahead order on alignment. The
normal time projection gives \(g\ge s\), while the
singleton \(e\ge1\) persists. The same zero-velocity test
then has gain at least \(1+2s\). If the \(B\)-marginal
vanishes, \(G\) depends on \(C\) alone and has global
slope \(P\), by \eqref{eq:source-28}. Thus the actual
profile satisfies \(p\ge P>u_*/2\), and the high
longitudinal comparison contradicts calibration.

It remains to treat \(a=0\), so \(c=b>0\).
Put the static coordinate \(C\) in the nonnegative group.
The external slices now have strict mixed slack also for
\(B\le C+R\), and the pure domain in \(B,R\) is full.
Prediction \eqref{eq:source-12} gives zero pure \(R\)-rate
on \(B>R\). The two-variable pure order rule therefore gives
\begin{equation}\label{asp:decreasing-zero-pure}
G(B,R)=Q_2 B\qquad(B\ge R),
\end{equation}
with \(Q_2\ge0\); the equality at \(B=R\) follows by
continuity. In particular this applies on alignment when
\(C\ge0\). Both negative coordinates can be shifted to
equal finer tied cutoffs in the upper envelope; the other mixed
constraints have large slack after that shift.

Suppose \(Q_2\le u_*/2\). Test the upper envelope
backwards from zero with \(v_2=1\) and
\(v_1\in[\rho,1-\rho]\). At interior negative times
\(C>0\), while \(A\) is fixed. The pure backward cost
per unit time is \(Q_2\), by
\eqref{asp:decreasing-zero-pure}. With \(B,R\) frozen
at their finer tied shift, the remaining scalar coordinates are
\[
X=C,\qquad y=Y,\qquad Z=D,
\]
with \(A\) fixed. Because \(C>0\), the first shear
constraint has slack; the only binding local constraint is
\(D\le C+Y\). The direct longitudinal velocity-time
projection gives \(e\ge f_s(p)\), and the time rate
is at least \(s\). The fixed normal measurements rebase
with bounded ambiguity by \eqref{eq:source-10}, with
\(B,R\) fine. Thus \eqref{eq:source-20} holds with
\(t_0=s\), also in the path tangents. The scalar comparison
gives \(s+h-\kappa\); the additional normal contribution is
\[
k_B-Q_2\ge1-\kappa-u_*/2=h-\kappa.
\]
The total is at least \(1+2s-2\kappa\), before an
arbitrarily small clipping and control loss, contradicting
\eqref{asp:contradiction}.

Finally suppose \(Q_2>u_*/2\). Use actual paths
backwards from zero with \(v_2=0\) and
\(v_1=m\in[\rho,1-\rho]\). At negative times \(C>0\),
so the pure bound gives \(n_1\ge Q_2\) locally and
therefore \(g\ge h\) at generic aligned states and in
their inherited tangents. Along each path, \(C\) alone
decreases; \(B\) is static, while \(R,Y,D,A\) strictly
increase. Apply the signed rule at almost every point of that
path, with pure \(C\)-slope \(p_0\). The required
overlap and slack hold: increase \(R,Y\) substantially
more than the shifts of the other nonnegative coordinates.

If \(p_0>u_*/2\), use \(v_1=v_2=0\) in this
tangent. The pure and projection bounds give \(e,g\ge h\),
and \(t\ge s\), yielding \(1+2s\).
If \(p_0\le u_*/2\), use \(v_1=1,v_2=0\).
We verify exactness of \(p=p_0\) where the gain is
transferred. The centered comparison ray has
\[
C=-y,\qquad B=0,\qquad R=y,\qquad y<0.
\]
Perturb to \(B>0\), \(R=B-C\), with \(C\)
near \(-y\). Away from \(y=0\), choose
\[
\beta\in(-C/m,(B-C)/m),\qquad \beta<y.
\]
This is possible for sufficiently small positive \(B\),
since at \(B=0,C=-y\) both endpoints tend to
\(y/m<y\). Then \(C>-m\beta\), and the
remaining cutoffs are strictly finer than their path values
\[
0,\quad m\beta,\quad\beta,\quad(1-m)\beta,\quad\beta
\]
for \(B,R,Y,D,A\), respectively. In particular
\(R=B-C>m\beta\); \(D=C+y\) is near zero
and exceeds \((1-m)\beta<0\); and
\(A=B+y>\beta\). Overlap therefore gives exact
\(p=p_0\) at generic nearby aligned states.
Integrate on these perturbed segments and pass to the ray.
Discarding the nonnegative \(r\)-term gives gain at least
\[
s+k_C-p_0+h\ge1+2s+\kappa.
\]
Clip endpoint velocities into the prescribed box and use
Lemma~\ref{ent:control}. This proves the contradiction in
the last decreasing-aspect case.
\end{proof}

\subsection{Conclusion of the geometric estimate}

\begin{proof}[Completion of the proof of Theorem~\ref{geo:main}]
If the claimed exponent bound failed, the calibrated construction would
give the profile used at the start of this section.
Proposition~\ref{asp:stationary} excludes \(\lambda=0\).
Propositions~\ref{asp:increasing-positive} and
\ref{asp:increasing-zero} exclude \(\lambda>0\), and
Proposition~\ref{asp:decreasing} excludes \(\lambda<0\).
These exhaust \(a\ge0,b\le1,\lambda=a-b\).

Each comparison is in a fixed tangent problem. Fix the finite parameters
and any large auxiliary velocity bound first. For an upper-envelope
comparison, then choose \(H_0\) sufficiently large for the resulting
compact test region. The admissible depth perturbations and velocity
clipping parameters may subsequently depend on these choices.
The scalar losses are at most \(2\kappa\) in the uses above,
and the local constrained comparison reserves less than \(4\kappa\).
Choose the additional endpoint, clipping, and control losses so that,
together with the scalar or local constrained loss already incurred,
the total in the selected comparison is strictly below \(10\kappa\).
These are alternative comparisons, so their losses are not summed
over the mutually exclusive cases. Static coordinates at \(a=1\)
or \(b=1\) belong to the nonnegative group, as permitted in
Lemma~\ref{ent:signed-envelope}; no division by a zero negative
speed is used. When \(s=1\), the branches requiring
\(0<P\le u_*/2\) are empty, and the zero-rate and high-rate
comparisons cover the remaining possibilities.
The finite realizing grids and large-scale limit are chosen after
these finite tests, as in Lemma~\ref{ent:tangent-realization}.
The strict inequality \eqref{asp:contradiction} is therefore
impossible, proving \eqref{eq:source-4}.
\end{proof}
